\documentclass[11pt]{article}

\usepackage{color,hyperref,amsmath}
\usepackage[sc,osf]{mathpazo} % Palatino font, w/ smallcaps & old-style figures
\usepackage[a4paper, total={5in, 8.25in}]{geometry}

\usepackage[utf8]{inputenc}
\usepackage{graphicx}
\usepackage{longtable}
\usepackage{enumitem}
\usepackage{tikz}
\usetikzlibrary{shadows.blur, positioning}

\newcommand{\code}[1]{\texttt{#1}}
\newcommand{\cmd}[1]{\texttt{#1}}

% The 5in measure is narrow for unbreakable \code identifiers; let TeX
% stretch interword space rather than overrun the margin.
\setlength{\emergencystretch}{1.5em}

% Default \tabcolsep (6pt each side of every column) adds up fast in
% this narrow 5in measure -- e.g. a 4-column table pays 48pt (0.67in)
% of padding alone. Trim it once, globally, rather than fighting
% column-width arithmetic per table.
\setlength{\tabcolsep}{4pt}

% Simple boxed CAUTION callout, styled plainly (no tcolorbox dependency).
\newcommand{\caution}[1]{%
  \par\noindent
  \fbox{\parbox{4.7in}{\textbf{CAUTION:} #1}}
  \par\smallskip
}

\title{WHAM-PFMG474-V4\\ \large Firmware Build, Flash, and Verification \\ Standard Operating Procedure}
\author{}
\date{August 2026}

\begin{document}
\maketitle

\begin{center}
\begin{tabular}{p{0.45in} p{0.75in} p{0.5in} p{2.7in}}
  \hline
  \textbf{Rev} & \textbf{Date} & \textbf{Author} & \textbf{Description} \\
  \hline
  A & 2026-08-31 & -- & Initial issue: build, bootstrap flash, serial
    reflash. No power-stage procedures yet (Section~\ref{sec:scope}). \\
  \hline
\end{tabular}
\end{center}

\section{Purpose}

This document specifies the procedure for building, flashing, and
verifying WHAM-PFMG474-V4 firmware, and for confirming a board is
running the expected firmware version before it is put into service or
handed off.

\section{Scope}
\label{sec:scope}

Covers firmware build and flashing only. \textbf{This firmware currently
implements no PWM, HRTIM, or power-stage control} -- it identifies
itself over serial and accepts remote firmware updates, nothing more.
This SOP will need a companion procedure for power-stage bring-up,
calibration, and safety interlocks once that functionality exists; do
not treat this document as covering any high-voltage or power-stage
operation.

\section{Prerequisites}

\subsection{Equipment}

\begin{center}
\begin{tabular}{p{1.5in} p{3.2in}}
  \hline
  \textbf{Item} & \textbf{Notes} \\
  \hline
  WHAM-PFMG474-V4 board & Powered per its own hardware documentation
    (not covered here). \\
  ST-Link (or equivalent SWD debug probe) & Required for the bootstrap
    flash (Section~\ref{sec:bootstrap}) and for recovery if a board
    stops responding over serial entirely. \\
  USB-to-serial adapter & Wired to USART2: adapter TX~$\to$~PA3
    (USART2\_RX), adapter RX~$\to$~PA2 (USART2\_TX), GND~$\leftrightarrow$~GND. \\
  \hline
\end{tabular}
\end{center}

\noindent\textit{(TX\slash RX wiring above is crossed: adapter TX goes
to the MCU's RX pin and vice versa, as usual for a point-to-point serial
link.)}

\subsection{Software}

\begin{center}
\begin{tabular}{p{1.5in} p{3.2in}}
  \hline
  \textbf{Tool} & \textbf{Notes} \\
  \hline
  STM32CubeIDE & Project build and the bootstrap SWD flash. \\
  \code{stm32flash} & AN3155 USART bootloader client. macOS:
    \code{brew install stm32flash}. \\
  Python 3 + \code{pyserial} & Runs \code{python/wham\_serial\_flash.py}.
    \code{pip install pyserial}. \\
  \hline
\end{tabular}
\end{center}

\caution{This firmware has not yet been through an ESD-controlled
production process. Use standard ESD precautions (grounding strap,
static-safe surface) when handling the bare board.}

\section{Procedure: Build the Firmware}
\label{sec:build}

\begin{enumerate}[label=\arabic*.]
  \item Open the project in STM32CubeIDE and press \textbf{Refresh}
    (\code{F5}) on the project if any files were added outside the IDE
    since it was last opened.
  \item Build the \textbf{Debug} configuration (\code{Project~$\to$~Build
    Project}, or from the command line: \code{cd Debug \&\& make all -j4}).
  \item Confirm the build produced no warnings. A clean build produces
    \code{WHAM-PFMG474-V4.elf} in \code{Debug/}.
  \item \textbf{Generate the \code{.bin}.} This project's build
    configuration does not have the ``Convert to binary file'' post-build
    step enabled, so it must be generated by hand:
\begin{quote}\ttfamily\small\raggedright
arm-none-eabi-objcopy -O binary \textbackslash\\
\hspace*{1em}Debug/WHAM-PFMG474-V4.elf \textbackslash\\
\hspace*{1em}Debug/WHAM-PFMG474-V4.bin
\end{quote}
\end{enumerate}

\section{Procedure: Bootstrap Flash (SWD)}
\label{sec:bootstrap}

Required the \emph{first} time a board receives firmware with the
serial-bootloader feature (\code{BOOT} command), since that feature is
what the serial reflash path in Section~\ref{sec:serial-flash} depends
on -- a board cannot use it to receive the very first firmware that adds
it. Also use this procedure to recover a board that no longer responds
to \code{BOOT} over serial at all.

\begin{enumerate}[label=\arabic*.]
  \item Connect the ST-Link to the board's SWD header and to the host
    computer.
  \item In STM32CubeIDE, press \textbf{Run} (or \textbf{Debug}) to build
    and flash via the ST-Link, as normal for this IDE.
  \item Proceed to Section~\ref{sec:verify} to confirm the board is
    running the expected firmware.
\end{enumerate}

\section{Procedure: Serial Reflash}
\label{sec:serial-flash}

Use this for all firmware updates \emph{after} the bootstrap flash in
Section~\ref{sec:bootstrap} has been done once. This is the
remote-friendly path: no physical BOOT0\slash NRST\slash SWD access is
required, only the USART2 serial link.

\begin{enumerate}[label=\arabic*.]
  \item With the board powered and running its current firmware, and the
    USB-serial adapter connected, run:
\begin{quote}\ttfamily\small\raggedright
python3 python/wham\_serial\_flash.py \textbackslash\\
\hspace*{1em}--port /dev/cu.usbserial-XXXXX
\end{quote}
    substituting the actual serial device path. (On macOS, a single
    USB-serial adapter shows up as both a \code{cu.*} and \code{tty.*}
    device node; the script sees this as two candidates and will not
    auto-detect -- \code{--port} must be passed explicitly. Prefer the
    \code{cu.*} node.)
  \item The script will: send \code{BOOT} and wait for the
    acknowledgement; run \code{stm32flash} to write and verify
    \code{Debug/WHAM-PFMG474-V4.bin}; then start the new firmware
    automatically (no manual reset needed -- see
    Section~\ref{sec:reset-note}).
  \item Confirm the script reports \texttt{[done]} with firmware
    written and verified before proceeding.
  \item Proceed to Section~\ref{sec:verify}.
\end{enumerate}

\subsection{Note on automatic restart after flashing}
\label{sec:reset-note}

Early testing of this procedure found that the newly-flashed firmware
could come up unresponsive over serial until the board was manually
power-cycled, because the STM32 ROM bootloader's \emph{Go} command (used
to start the new firmware without a physical reset) does not perform a
true chip reset, and left the processor's vector table pointer aimed at
the bootloader's own vector table rather than the application's. This
has been fixed in firmware (\code{boot\_jump.c}) -- a manual reset should
not be necessary. If a board flashed under this SOP does not respond to
the verification step below without a manual reset, that is a
regression and should be reported, not treated as expected behavior.

\section{Procedure: Verification}
\label{sec:verify}

\begin{enumerate}[label=\arabic*.]
  \item Open a serial connection to the board at \textbf{9600 8N1}.
  \item Send \code{*IDN?} followed by a line ending (\code{\textbackslash
    r}, \code{\textbackslash n}, or both).
  \item Confirm a reply of the form:
    \begin{quote}
      \code{OK <board name> <hardware rev> <firmware version>}
    \end{quote}
    e.g.\ \code{OK WHAM-PFMG474-V4 REVA v0.6}.
  \item Confirm the reported firmware version matches the version that
    was just built\slash flashed.
\end{enumerate}

Acceptance criterion for this SOP: Step 3 above returns a well-formed
reply with the expected version, on the first query, with no manual
reset performed since the flash in Section~\ref{sec:build} or
\ref{sec:serial-flash}.

\section{Troubleshooting}

\begin{center}
\begin{tabular}{p{1.7in} p{3in}}
  \hline
  \textbf{Symptom} & \textbf{Cause / action} \\
  \hline
  \code{stm32flash} reports \code{Failed to init device, timeout.}
    before any write occurs &
    Not a firmware issue -- this probe doesn't depend on what's running
    on the chip. Check board power, serial wiring (TX\slash RX not
    crossed, common GND), and that no other program has the port
    open. \\
  \code{BOOT} gets no reply, but power/wiring check out &
    The firmware currently on the board may predate the \code{BOOT}
    command (run the bootstrap flash, Section~\ref{sec:bootstrap}) or
    may be a board affected by the \code{Go}-jump issue in
    Section~\ref{sec:reset-note} -- power-cycle once, then retry. \\
  Verification (Section~\ref{sec:verify}) fails after a serial reflash &
    Power-cycle the board once and repeat verification. If it then
    succeeds, this is the regression described in
    Section~\ref{sec:reset-note} -- report it. If it still fails, treat
    as a failed flash and repeat Section~\ref{sec:serial-flash} from the
    start. \\
  \hline
\end{tabular}
\end{center}

\section{References}

\begin{itemize}
  \item \code{AGENTS.md} -- project orientation and current functionality.
  \item \code{docs/command\_reference.md} -- full serial command protocol.
  \item \code{docs/serial\_reflash\_guide.md} -- detailed reflash
    mechanism, including the \code{Go}-jump bug referenced in
    Section~\ref{sec:reset-note}.
  \item \code{docs/changelog.txt} -- dated development history.
\end{itemize}

\end{document}
